\chapter{Python Environment, Data, and Repository Standard}
\label{ch:python}

\chapterguide%
  {set up the course project, install Python dependencies, and obtain the required datasets.}
  {run every laboratory reproducibly in Jupyter and maintain one professional GitHub repository for the course.}

\begin{keyidea}
This course uses \textbf{Python only}. All assessed experiments are provided as Jupyter notebooks and all reusable course code is written in Python.
\end{keyidea}

\section{Standard project structure}

Use one repository for the whole course. Keep code, notebooks, datasets, and generated evidence in predictable locations:

\begin{lstlisting}[style=mlterminal]
UTP-ARTIFICIAL-INTELLIGENCE/
  all_datasets/
    lab01-03_airports.csv
    lab01-03_routes.csv
    lab04-AirQuality.csv
    lab05-sensor_readings_24.csv
    lab06-ocr-handwritten-prescription.zip
    lab07-Tweets.csv
    lab08-insurance.csv

  all_experiments/
    lab01_search_foundations_worked.ipynb
    lab02_informed_search_worked.ipynb
    lab03_search_engineering_worked.ipynb
    lab04_fuzzy_anfis_worked.ipynb
    lab05_robotics_sensor_control_worked.ipynb
    lab06_handwriting_ocr_worked.ipynb
    lab07_airline_sentiment_worked.ipynb
    lab08_genetic_algorithm_regression_worked.ipynb

  src/search_toolkit/
  tests/
  config/
  results/lab01/ ... results/lab08/
  requirements.txt
  run_all_notebooks.py
  README.md
\end{lstlisting}

\begin{keyidea}
\textbf{Do not rename supplied files casually.} The notebooks use the filenames above so that a clean checkout can be reproduced without editing paths in multiple places.
\end{keyidea}

\section{Dataset acquisition}

Download each dataset manually from its Kaggle webpage. No Kaggle API key or token is required. Sign-in may be required by Kaggle. The course project already uses the local filenames shown below.

\begin{mltable}
\footnotesize
\renewcommand{\arraystretch}{1.08}
\begin{tabularx}{\linewidth}{@{}p{0.36\linewidth}p{0.12\linewidth}X@{}}
\toprule
\rowcolor{MLPaleBlue!92}
\textbf{Kaggle dataset} & \textbf{Lab(s)} & \textbf{Local file and purpose} \\
\midrule
\href{https://www.kaggle.com/datasets/open-flights/flight-route-database}{\textbf{OpenFlights Flight Route Database}}
& 1--3 & \texttt{lab01-03\_routes.csv}; directed airline routes for uninformed and informed search. \\
\addlinespace
\href{https://www.kaggle.com/datasets/open-flights/airports-train-stations-and-ferry-terminals/data}{\textbf{OpenFlights Airports, Train Stations and Ferry Terminals}}
& 1--3 & \texttt{lab01-03\_airports.csv}; IATA codes and coordinates for graph construction and Haversine heuristics. \\
\addlinespace
\href{https://www.kaggle.com/datasets/fedesoriano/air-quality-data-set}{\textbf{Air Quality Dataset}}
& 4 & \texttt{lab04-AirQuality.csv}; fuzzy control and ANFIS regression. \\
\addlinespace
\href{https://www.kaggle.com/datasets/uciml/wall-following-robot}{\textbf{Wall-Following Robot Sensor Readings}}
& 5 & \texttt{lab05-sensor\_readings\_24.csv}; 24 ultrasonic sensor readings plus movement class. \\
\addlinespace
\href{https://www.kaggle.com/datasets/mrdude20/doctor-handwriting-recognition-dataset}{\textbf{Doctor Handwriting Recognition Dataset}}
& 6 & \texttt{lab06-ocr-handwritten-prescription.zip}; 89 labelled handwriting images for OCR evaluation. \\
\addlinespace
\href{https://www.kaggle.com/datasets/chandana890/twitter-airline-sentiment-dataset}{\textbf{Twitter Airline Sentiment Dataset}}
& 7 & \texttt{lab07-Tweets.csv}; a single CSV for three-class sentiment classification. \\
\addlinespace
\href{https://www.kaggle.com/datasets/mirichoi0218/insurance}{\textbf{Medical Cost Personal Dataset}}
& 8 & \texttt{lab08-insurance.csv}; regression data for genetic-algorithm optimisation. \\
\bottomrule
\end{tabularx}
\end{mltable}

\subsection{Download procedure}

\begin{mlsteps}
  \item Open the Kaggle dataset page and choose \textbf{Download}.
  \item Extract the archive when required and copy only the files used by the laboratory into \texttt{all\_datasets/}.
  \item Keep the filenames exactly as listed in the table.
  \item Run the corresponding notebook from the project root and confirm that its data-validation cell passes before continuing.
\end{mlsteps}

\begin{pitfall}
For \texttt{FileNotFoundError}, check the current working directory and filename first. Correct the path rather than copying duplicate datasets into the project root.
\end{pitfall}

\section{Install Python and create the environment}

Install a current Python 3 release from the \href{https://www.python.org/downloads/}{official Python downloads page}. Then create a project-specific virtual environment from the repository root:

\begin{lstlisting}[style=mlterminal]
# Windows
py -m venv .venv
.venv\Scripts\Activate.ps1
py -m pip install --upgrade pip
py -m pip install -r requirements.txt

# macOS / Linux
python3 -m venv .venv
source .venv/bin/activate
python3 -m pip install --upgrade pip
python3 -m pip install -r requirements.txt
\end{lstlisting}

The supplied \texttt{requirements.txt} includes NumPy, pandas, SciPy, scikit-learn, matplotlib, Jupyter, NetworkX, OpenCV, pytesseract, and pytest. Lab 6 also requires the system-level Tesseract OCR executable.

\begin{lstlisting}[style=mlterminal]
python --version
python -m pip --version
jupyter lab

# Lab 6 only
tesseract --version
\end{lstlisting}

\begin{keyidea}
Create the environment once, then reactivate it at the start of each session. Install packages into the project environment rather than a shared global Python installation.
\end{keyidea}

\section{Run and verify the course project}

Open notebooks from \texttt{all\_experiments/} in numerical order. Each worked notebook is designed to run from top to bottom and write evidence to its matching \texttt{results/labXX/} folder.

Run the automated search-toolkit tests with:

\begin{lstlisting}[style=mlterminal]
# macOS / Linux
PYTHONPATH=src pytest -q

# Windows PowerShell
$env:PYTHONPATH = "src"
pytest -q
\end{lstlisting}

To execute the complete notebook suite from a clean project checkout:

\begin{lstlisting}[style=mlterminal]
python run_all_notebooks.py
\end{lstlisting}

\section{GitHub repository standard}

Every student maintains one course repository. Use meaningful commits and tag the completed state of each laboratory.

\begin{mltable}
\begin{tabularx}{\linewidth}{@{}p{0.25\linewidth}X@{}}
\toprule
\rowcolor{MLPaleBlue!92}
\textbf{Item} & \textbf{Required standard} \\
\midrule
Repository & One repository for Labs 1--8; clear root-level \texttt{README.md}. \\
Notebook & Executed top-to-bottom with outputs visible and no unresolved exceptions. \\
Source code & Reusable logic moved to \texttt{src/} when appropriate; avoid large duplicated code blocks. \\
Results & Save required CSV/PNG evidence under \texttt{results/labXX/}. \\
Datasets & Do not commit large archives or unnecessary duplicates. Follow the course \texttt{.gitignore}. \\
Commits & Use concise messages such as \texttt{lab05: add noise robustness evaluation}. \\
Tags & Create \texttt{lab01-final}, \ldots, \texttt{lab08-final} after verification. \\
\bottomrule
\end{tabularx}
\end{mltable}

A typical submission sequence is:

\begin{lstlisting}[style=mlterminal]
git status
git add all_experiments results src tests README.md
git commit -m "lab07: complete sentiment evaluation"
git push
git tag lab07-final
git push origin lab07-final
\end{lstlisting}

\begin{pitfall}
Do not use GitHub as a screenshot archive. Commit the notebook, source code, concise machine-readable results, and only the figures needed to support the analysis.
\end{pitfall}

\section{Professional laboratory workflow}

Use the same six-stage workflow in every lab:

\begin{enumerate}
  \item \textbf{Validate the data.} Check shape, schema, missing values, labels, and assumptions.
  \item \textbf{Establish a baseline.} Define what ``better'' means before tuning a model.
  \item \textbf{Implement the method.} Keep parameters explicit and code readable.
  \item \textbf{Run a controlled experiment.} Use fixed splits/seeds and change one design decision at a time.
  \item \textbf{Evaluate and interpret.} Report quantitative metrics, inspect failures, and explain limitations.
  \item \textbf{Reproduce and publish.} Restart the notebook, run all cells, save evidence, commit, push, and tag.
\end{enumerate}

\begin{keyidea}
A successful laboratory is not ``code that ran once''. It is a reproducible experiment with a stated question, controlled procedure, measurable evidence, and a defensible interpretation.
\end{keyidea}
